% Folder 91, batch 2 (archivists' pages 21-40).
% Pass: Opus 5.5 (claude-opus-5-5), 2026-10-03 — first pass, unchecked against the pages by a human.
%
% Pages 29 and 32 are blank versos (ink show-through only) and are skipped.
% Pages 30-31 are a typed letter to Anantharaman dated 11.9.1967. Notation
% fixed for the batch: \operatorname{Hom}, \prod_{X/S} P/X,
% f_* for Weil restriction, \check{\mathbb{V}}(\mathcal{F}).
\documentclass[11pt,a4paper]{article}
\input{../preamble/grothendieck}

\folder{91}
\batch{2}
\pages{21}{40}
\foldertitle{Autour de Néron / Greenberg-Néron. Foncteurs Hom (méthodes non-projectives) : notes manuscrites (s.d.), lettre (1967).}
\dating{[à partir de 1964-vers 1970]}
\watermark{Édition de démonstration}

\begin{document}

\page{21}
\note{le haut de la page achève un argument commencé avant ce lot.}
\textbf{Pb} Simplifier les \add{\uncertain{démonstrations}} raisonnements, en éliminant
les hypothèses de propreté sur $X$, $Z$ et le \ill{} résultat délicat de
l'existence de $M_{\ill{}}$, $M_{\ill{}}$, en imposant l'hypothèse de
\uncertain{régularité} \uncertain{suivante}\marginal{?} :
\[
\operatorname{Hom}\bigl(\omega^{\circ}(\Omega^{1}_{Y_0/S_0}), \mathcal{O}_{X_0}\otimes G\bigr)
\longrightarrow
\operatorname{Hom}\bigl(\omega^{\circ}(\Omega^{1}_{X_0/S_0}), \mathcal{O}_{Z_0}\otimes G\bigr)
\]
injectif pour tout $G$ \uncertain{quasi-cohérent} sur $S_0$.

\note{toute la suite de la page est biffée de deux longs traits obliques ; on
en donne ce qui se lit. En marge, au crayon et entouré : « SGA 3 III 4.2 ».}

\struck{\textbf{Prop.}} \struck{Le \add{\uncertain{critère}} \ill{} géométrique, on a \ill{} : ceci.
$S$ étant \ill{} \uncertain{local} \ill{} \uncertain{point générique} $y$,
$S_0$ \ill{} d'un idéal $I$ de carré nul, \ill{} \ill{} (\ill{}),
\uncertain{disons} $S_0 = V(J)$, et $u_0 : X_0 \to Y_0$ avec $v_0 = u_0\varphi_0$,
\ill{} \ill{} \ill{} \uncertain{pour que} \ill{}
cette condition, \uncertain{prenons} \ill{} \ill{} d'un \uncertain{voisinage}
\ill{} \uncertain{l'obstruction} \ill{} \ill{}
$S'_{0_s} \supset S_{0_s}$ donné, \ill{}
\uncertain{condition} (\ill{} \ill{}) $v$ étant \ill{} la \uncertain{question} \ill{}
\ill{} immédiate, \ill{} \uncertain{obstruction}}
\struck{[cf. $\operatorname{Ext}^{1}_{\mathcal{O}_{P_0}}(P_0;
\underline{I}_{\Gamma_0}, \mathcal{O}_{\Gamma_0}\otimes_{S_0} J)$]}

\page{22}
$S$ \add{\uncertain{unité}} \uncertain{muni} d'un idéal \add{\uncertain{loc. nilp.}} de carré nul $J$,
$Z \xrightarrow{v} X, Y$ sur $S$ \marginal{d'où $v_0 : Z_0 \to X_0, Y_0$ sur $S_0$}
\note{le petit ajout entouré « d'où $v_0 : Z_0 \to X_0, Y_0$ sur $S_0$ » est écrit
en haut à droite, sous $Z \to X, Y$.}

$S_0 = V(J)$,

$u_0 : X_0 \to Y_0$ tel que $u_0\varphi_0 = v_0$.

On va étudier l'obstruction à relever $u_0$ en $u : X \to Y$
[tel que $u\varphi = v$]. À priori, \struck{l'obstruction à \uncertain{relever}}
\add{les relèvements de} $u_0$ \ill{} \ill{} \ill{} \uncertain{comme} \ill{}, condition
supplémentaire $u\varphi = v$) \ill{}, correspondant aux splittings d'une extension sur $P_0$
\marginal{cf. l'obstruction \ill{} \uncertain{ce que} \ill{} \uncertain{revient} à l'idée d'un
$u : J \to G$, \ill{} \uncertain{puis localisation}\,…}
\[
0 \to \mathcal{O}_{\Gamma_0}\otimes_{S_0} J \to E \to \underline{I}_{\Gamma_0} \to 0 .
\]
\marginal{SGA 3 III 4.2}
\note{la référence « SGA 3 III 4.2 », au crayon et entourée, est reliée par un
trait à « les relèvements » de la ligne suivante.}
Les relèvements \uncertain{fixant} $v_0$ correspondent \ill{} aux splittings
d'une extension \ill{} sur $P'_0 = Z_0 \times_{S_0} Y_0$
\[
0 \to \mathcal{O}_{\Gamma'_0}\otimes_{S_0} J \to E' \to \underline{I}_{\Gamma'_0} \to 0 .
\]
D'ailleurs, on a \uncertain{évidemment} \uncertain{un} homomorphisme
d'extensions, \uncertain{compatible} \ill{} $P'_0 \to P_0$

\begin{tikzcd}[column sep=small]
0 \arrow[r] & \mathcal{O}_{\Gamma_0}\otimes_{S_0} J \arrow[r] \arrow[d] & E \arrow[r] \arrow[d] & \underline{I}_{\Gamma_0} \arrow[r] \arrow[d] & 0 \\
0 \arrow[r] & \mathcal{O}_{\Gamma'_0}\otimes_{S_0} J \arrow[r] & E' \arrow[r] & \underline{I}_{\Gamma'_0} \arrow[r] & 0
\end{tikzcd}

(en fait, $E'$ se déduit de $E$ par extension de la base
\uncertain{par} $P'_0 \to P_0$ \ill{} \ill{}

\begin{tikzcd}[row sep=small]
 & P_0 \arrow[d] \\
Z \arrow[r] & X
\end{tikzcd}

; il \uncertain{faut} \uncertain{préciser} \ill{} $\varphi$ \ill{}].
\struck{\uncertain{Conditions}, \uncertain{puisque} \ill{}}
\add{D'ailleurs, la} \ill{} \uncertain{signifie} que
\uncertain{l'image} d'un splitting, \ill{} \uncertain{la relation}
\ill{} les splittings de la première
\uncertain{compatibles} \ill{} celui \uncertain{donné} pour la seconde.

\page{23}
Ceci fait, regardons, pour $G$ \uncertain{quasi-cohérent} variable
sur $S_0$, les foncteurs
$\underline{\operatorname{Ext}}^{1}(\ , \mathcal{O}_{\Gamma_0}\otimes_{S_0} G)$,
$q'_{*}\underline{\operatorname{Ext}}^{1}(\ , \mathcal{O}_{\Gamma'_0}\otimes_{S_0} G)$
relatifs aux deux lignes, \uncertain{on} \uncertain{trouve} \ill{}
\add{(via les théorèmes de \uncertain{représentabilité})}
\uncertain{un} \uncertain{diagramme} \uncertain{commutatif} de \uncertain{suites} \uncertain{exactes} (D)

\begin{tikzcd}
 & L \arrow[r, "i"] & M \arrow[r, "j"] & N \arrow[r] & 0 \\
0 \arrow[r] & L' \arrow[r, "i'"] \arrow[u, leftrightarrow, "\alpha"'] & M' \arrow[r, "j'"] \arrow[u, leftrightarrow, "\beta"'] & N' \arrow[r] \arrow[u, "\gamma"'] & 0
\end{tikzcd}

\note{les flèches verticales $\alpha$ et $\beta$ portent une pointe à chaque
bout : une pointe vers le bas semble surchargée d'une pointe vers le haut ;
$\gamma$ monte. On garde les deux pointes.}

de plus, \ill{} \ill{} \uncertain{un} \uncertain{homomorphisme} can.\ $\sigma$

\begin{tikzcd}[row sep=small]
 & & L \arrow[dll, "\sigma"'] \\
J \arrow[d, dashed, "\lambda"'] & & L' \arrow[u, "\alpha"'] \\
G & &
\end{tikzcd}

\note{$J$, $L$, $L'$ et la flèche $\sigma$ sont entourés d'un même trait ; la
flèche pointillée $J \to G$ est ajoutée en dessous, et la lettre près de $J$
(lue $\lambda$) est surchargée.}

correspondant à l'identité
$\mathcal{O}_{\Gamma_0}\otimes_{S_0} J \to \mathcal{O}_{\Gamma_0}\otimes_{S_0} J$
($J$ \add{$= G$}), et un splitting \add{$M' \xrightarrow{\pi'} L'$} de la deuxième ligne

\[
L' \xleftarrow{\ \pi'\ } M' .
\]

Ceci étant, pour un $\lambda : J \to G$ donné
\marginal{\ill{} \ill{}},
les relations \ill{} \uncertain{i.e.}\ \uncertain{tels} \uncertain{que} les
\uncertain{relations} \ill{} \ill{} \ill{}, les homom.\ $M \xrightarrow{\mu} G$
qui \add{relèvent} \ill{} i.e.\ $\boxed{\mu i = \lambda\sigma}$
\ill{} $L \xrightarrow{\lambda\sigma} G$, et tels que
\[
\mu\beta = (\lambda\sigma\alpha)\pi' .
\]
\struck{\ill{}} Pour \ill{} qu'il \uncertain{en} existe un,
\marginal{condition de non ramification \uncertain{de} $u_0$}
il f.\ et s.\ que
\begin{enumerate}
\item[1°)] $(\lambda\sigma)(\operatorname{Ker} i) = 0$ \quad i.e.\
$\lambda(R_1) = 0$, \quad où $R_1 = \sigma(\operatorname{Ker} i)$
\marginal{Alors on peut remplacer $L$ par $L/\operatorname{Ker} i$}

\struck{[i.e.\ $\lambda$ se factorise \ill{} \ill{} \ill{}
$\bar\lambda : \bar J = J/R_1 \to G$]}

\item[2°)] [\uncertain{Tenant} \uncertain{compte} que \struck{\ill{}}
\uncertain{$\gamma$ est injectif}]
$\lambda\sigma\alpha\pi'$ \uncertain{nulle} sur $\operatorname{Ker}\beta$,
i.e.\ $\lambda(R_2) = 0$, où $R_2 = \sigma\alpha\pi'(\operatorname{Ker}\beta)$.
\end{enumerate}
\note{la suite (une éventuelle troisième condition) passe à la page suivante.}

\page{24}
\begin{tikzcd}[column sep=small, row sep=small]
\operatorname{Ker} i \arrow[rr, dashed] \arrow[d] & & L \arrow[dll, dashed] & \\
J & & L' \arrow[ll, dashed] \arrow[u, dashed, "?"'] & M' \arrow[l, dashed] \\
 & & & \operatorname{Ker}\beta \arrow[u, dashed] \arrow[ulll, bend left=10]
\end{tikzcd}

\note{l'étiquette de la flèche pointillée $L' \to L$ est peu lisible (« 2 » ?).}

\note{le diagramme suivant est encadré, à droite du précédent.}

\begin{tikzcd}
L + M' \arrow[r, two heads, "{(i,\beta)}"] \arrow[d, "{(\sigma,\, \sigma\alpha\pi') = \rho}"'] & M \arrow[ddl, dashed, "\mu"] \\
J \arrow[d, "\lambda"'] & \\
G &
\end{tikzcd}

\marginal{\uncertain{Surj.} (entouré, sous $(i,\beta)$) : condition de non ramification.
Le foncteur en $G$ est repr.\ par $\operatorname{Coker}\rho$, dont la formation
commute à tout chgt de base.}

En résumé, \ill{} \uncertain{opération} \uncertain{qui} \uncertain{précède}
\ill{} $\lambda$ \uncertain{sur} $R = R_1 + R_2$,
\uncertain{l'annulation} de $\lambda$ \uncertain{sur} \ill{}
\uncertain{revient} à la factorisation de $\lambda$ par $J/R \to G$.
D'ailleurs, la formation des quotients \uncertain{ainsi} \uncertain{obtenus}
$J/R$ de $J$ commute \uncertain{à} \uncertain{l'extension} \uncertain{plate}
de la base $S$ [\uncertain{ou} \uncertain{générale} : \uncertain{à} \uncertain{la}
\uncertain{localisation}]. Ceci \uncertain{donne} la \uncertain{validité} des
\uncertain{conditions} \struck{7°} \add{7°} et 8°.
\note{« 7° » est récrit par-dessus un premier chiffre ; ces numéros renvoient
à une liste de conditions qui n'est pas dans ce lot.}

\note{ici une ligne biffée, illisible, puis un trait horizontal.}

\textbf{Corollaire} Soient $P/X/k$ sur $k$ \uncertain{alg.}\ clos\marginal{(\uncertain{déjà} \uncertain{non} \uncertain{trivial} !)},
$X$ propre sur $k$, $P/X$ loc.\ de \struck{présentation} \add{type} fini. Alors
$\prod_{X/k} P/X$ \uncertain{représentable} \uncertain{par} \uncertain{préschéma}
\uncertain{loc.} \uncertain{de} \uncertain{type} \uncertain{fini}.
\struck{\ill{}}

Si $P/X$ quasi-projectif \add{\uncertain{à} \uncertain{fibres} qu.-pr.}\ \uncertain{ou}
\uncertain{seulement} \uncertain{à} \uncertain{fibres} \uncertain{quasi-projectives},
alors $\prod_{X/k} P/k$ \ill{} essentiellement quasi-projectif
[i.e.\ \uncertain{ses} \uncertain{composantes} \uncertain{connexes} \uncertain{sont}
quasi-projectif sur $k$], \uncertain{dès} \uncertain{que} $\forall x \in$
\ill{} $\mathcal{O}_X$, $\exists y \in \bar{x}$ tel que
\add{\uncertain{il} \uncertain{suffit} que la fibre $P_y$ soit quasi-projective}.
\note{la fin de la page est surchargée : un ajout interlinéaire (« il suffit que
la fibre $P_y$ soit quasi-projective ») semble remplacer la fin de la phrase.}

\page{25}
Corollaire 1 : \uncertain{est} contenu dans le suivant :

\textbf{Théorème} $X$ propre et plat de prés.\ finie sur $S$ \struck{localement},
$P/X$ loc.\ de présentation finie sur $X$,
$Z_1, \ldots, Z_n$ des sous-préschémas \add{fermés} de $X$, plats et de
présentation finie sur $S$. On suppose
\begin{enumerate}
\item[(a)] Pour tout $i$, $\prod_{Z_i/S} P_{Z_i}/Z_i$ représentable
[p.\ ex.\ $Z_i/S$ radiciels, ou $Z_i/S$ et $P_{Z_i}/Z_i$ projectifs]
\item[(b)] Pour tout $i$, $X$ est normalement plat le long de $Z_i$
[p.\ ex.\ $Z_i$ \uncertain{simples}$/S$ et $X/S$ simple au \uncertain{voisinage} de $Z_i$]
\item[(c)] Pour tout $s \in S$, et tout $x \in$ \struck{Ass} $\mathcal{O}_{X_{\bar s}}$,
il existe $i$ tel que $Z_{i\bar s} \cap \bar{x} \neq \emptyset$, i.e.\
$(\bigcup Z_i)_{\bar s} \cap \bar{x} \neq \emptyset$
[$\bar s = \operatorname{Spec}$ \uncertain{$k(\bar s)$}]
\end{enumerate}
Sous ces conditions,
\begin{enumerate}
\item[(i)] $\prod_{X/S} P/X$ est représentable par un préschéma loc.\ de
\struck{type} \add{présentation} finie et \struck{séparé} sur $S$
\item[(ii)] \struck{Si de plus} Pour toute partie \add{\uncertain{ouverte}}
quasi-compacte $U$ de $\prod_{X/S} P/X$, \add{(\uncertain{qui} est quasi-séparé)}
\struck{\ill{} \ill{} \ill{}}, le morphisme
\[
U \longrightarrow \prod_{1\le i\le n} \text{\struck{\ill{}}} \prod_{Z_i/S} P_{Z_i}/Z_i
\]
est quasi-affine.
\item[(iii)] Par suite, si les $\prod_{Z_i/S} P_{Z_i}/Z_i$ sont
\add{sur $S$} \struck{\ill{}} quasi-projectifs, p.\ ex.\ $Z_i/S$ finis et
$P_{Z_i}/Z_i$ projectifs, alors toute partie quasi-compacte de
$\prod_{X/S} P/X$ est quasi-projective sur $S$ [de façon précise, il
\uncertain{existe} \uncertain{un} \uncertain{Module} inversible $L$ sur
$\prod_{X/S} P/X$ qui est \ill{} ample \ill{} \ill{}$/S$]
\end{enumerate}
\note{le mot ajouté au-dessus de « quasi-compacte » au (ii), lu « ouverte »,
et l'ajout entouré « quasi-séparé » sont peu sûrs ; une ligne biffée suit
dans l'interligne. Au bas de la page, quelques mots serrés, illisibles,
achèvent le crochet du (iii).}

\page{26}
\textbf{Pb} Les composantes connexes de $\prod_{X/k} P/X$ \struck{\ill{}} sont-elles
de type fini sur $k$ ?

\textbf{Corollaire 2} Supposons, \add{pour} $P/X/k$, $X$ propre sur $k$,
\struck{\ill{} \ill{}}, $P/X$ loc.\ de type fini, $k$ \uncertain{corps}
(\uncertain{pas} \uncertain{nécessairement} alg.\ clos) que pour tout $x_i \in \operatorname{Ass}\mathcal{O}_X$,
$\exists\, x'_i \in \bar{x}$ tels que $k(x'_i)/k$ radiciel.
\add{[plus généralement, tels que $\prod_{k(x'_i)/k} P_{x'_i}/k(x'_i)$ représentable]}
\uncertain{Alors} $\prod_{X/k} P/X$ \uncertain{représentable}. P.\ ex.\ si $X$
\struck{\uncertain{lisse}} \uncertain{réduit} \ill{} \add{\uncertain{irréductible} \uncertain{aux}}
\struck{\uncertain{irréductible}} \struck{\uncertain{géométriquement}}, et a un point $x$
\uncertain{rationnel} sur $k$ \struck{\ill{}}, $\prod_{X/k} P/X$ existe.
\marginal{Si les $P_{x_i}$ sont quasi-projectifs, alors $\prod_{X/k} P/X$ est
\uncertain{essentiellement} quasi-projectif.}

\textbf{Corollaire 3} $P/X/k$ \ill{} $X/k$ \uncertain{comme} \uncertain{dessus}. Alors
$\exists\, k'/k$ extension finie \uncertain{séparable}, telle que
$\prod_{X'/k'} P'/X'$ représentable.
\note{les deux derniers corollaires sont encadrés à droite d'un long crochet.}

\textbf{Corollaire 4} Si $P/X$ à fibres \uncertain{quasi-projectives}, et
$X/k$ propre \struck{\ill{}}, dans $\prod_{X/k} P/X$ \uncertain{existe} et est
\uncertain{ess.}\ qu.-proj.\ (\uncertain{Comparer} : Cor.\ 1) \struck{\ill{}}

\textbf{Corollaire 5} $P/X/S$, $X$ propre et plat sur $S$ \uncertain{noeth.}
\uncertain{intègre}, $P/X$ loc.\ de type fini. Alors $\exists$
\struck{\uncertain{voisinage}} \ill{} \ill{} $U \neq \emptyset$, et
\uncertain{morphisme} fini \uncertain{étale} surjectif $U' \to U$, tel que
$\prod_{X_{U'}/U'} P_{U'}/X_{U'}$ représentable.
Si $X$ \uncertain{géométriquement} \ill{} \ill{} \uncertain{la} \uncertain{condition}
\uncertain{énoncée} dans Cor.\ 2, \ill{} \uncertain{peut} \ill{} prendre $U' = U$
\ill{}

\note{la marge gauche porte, en travers, des notes serrées et en partie
biffées, qu'on ne lit que par fragments : « \ill{} \ill{} \ill{} $P/X$ \ill{} \ill{} ??? »
et « \ill{} $P/X$ \ill{} \ill{} simple ?? » ; dans la marge du haut, un petit
tableau de traits verticaux et de mots (« \uncertain{différents} \ill{} »)
reste illisible.}

\page{27}
\textbf{Corollaire 6} $P/X/S$, avec $X \to S$ propre plat :
fibres géométriques intègres. Supposons qu'il existe une section $\sigma$ de
$X/S$ \struck{et $P/X$ \ill{} \ill{} \ill{} \ill{}}
\uncertain{le} \uncertain{long} \uncertain{de} \uncertain{laquelle}
$X/S$ \ill{} \add{\uncertain{lisse}}. Alors $\prod_{X/S} P/X$
représentable. Si $P \times_X (S,\sigma)$ est quasi-projectif,
$\prod_{X/S} P/X$ est essentiellement quasi-projectif.
\note{le chiffre de « Corollaire 6 » est écrit par-dessus un « 5 ».}

\textbf{Corollaire 7} $P/X/S$, avec $X \to S$ simple propre
\struck{[\uncertain{Comme} \uncertain{dessus}, $P/X$ \ill{} \ill{} \ill{}
\uncertain{fibres} \uncertain{géométriques} \uncertain{intègres}]}
\add{\uncertain{Alors}} \ill{} $\prod_{X/S} P/X$ est localement
\add{représentable} pour la topologie étale.

\uncertain{Remarque} \uncertain{sur} la \uncertain{top.}\ étale pour 2 \uncertain{top.}\
\ill{} (\uncertain{fid.}\ plat \uncertain{quasi-fini}), le résultat \uncertain{demeure}
\ill{} \uncertain{valable} \uncertain{sans} \uncertain{supposer} $X \to S$ simple, mais
seulement propre, plat, à fibres \struck{\uncertain{géométriquement} \uncertain{intègres}}
\add{\uncertain{sans} \uncertain{cycles}} \uncertain{immergés}, \uncertain{par}
essentiellement la \uncertain{même} \uncertain{démonstration}.

\textbf{Corollaire 8} $X$ et $Y$ propres sur corps $k$.
Alors $\underline{\operatorname{Isom}}_k(X,Y)$ est représentable. En
particulier, $\underline{\operatorname{Aut}}_k(X)$ est représentable.

\page{28}
\textbf{Corollaire 9} Soit $X/S$ propre et plat sur $S$
noethérien intègre. Alors $\exists\, U$ \uncertain{non} vide,
\add{et $U' \to U$ fini étale surjectif}
\struck{et $U' \to U$ fini \ill{} \ill{} \ill{}} tel que
$\underline{\operatorname{Aut}}_S(X) \times_S U'$ soit représentable.

\marginal{cf.\ Raynaud}
Il est plausible qu'on puisse \uncertain{même}
prendre $U' = U$, \uncertain{c'est-à-dire} \uncertain{à} \uncertain{une} question
d'effectivité de descente\,….
\note{la note « cf.\ Raynaud » est au crayon, dans la marge gauche ; le reste
de la page est blanc.}

\page{30}
\note{pp.\ 30--31 : lettre dactylographiée, en anglais, à Anantharaman, datée à
la main « 11.9.1967 » ; on la transcrit telle qu'elle est tapée, fautes de
frappe comprises, et sans la signature, absente.}

\marginal{11.9.1967}

Dear Anantharaman,

Matsumura proved that if $X$ is proper over a field $k$, then
$\underline{\operatorname{Aut}}_{X/k}$ is representable by a group scheme
locally of finite type over $k$. I think I can systematize the key step of his
argument in the following way. Consider a scheme $S$, and a morphism
\[
\phi : Z \longrightarrow X
\]
of $S$-schémes which are proper, flat and of finite presentation. Let $Y$ be
locally of finite présentation and separated over $S$, then $\phi$ induces a
homomorphisme of functors $u \rightsquigarrow u\phi$ :
\[
\phi' : \underline{\operatorname{Hom}}_S(X,Y) \longrightarrow \underline{\operatorname{Hom}}_S(Z,Y) .
\]
Then one can define a subfunctor of $\underline{\operatorname{Hom}}_S(X,Y)$
where $\phi'$ is ``unramified'' in a rather obvious sense, and this turns out to
be an ``open subfunctor'', say $\underline{\operatorname{Hom}}_S(X,Y;\phi)$.
Now look at the induced homomorphism
\[
\underline{\operatorname{Hom}}_S(X,Y;\phi) \longrightarrow \underline{\operatorname{Hom}}_S(Z,Y) .
\]
Using the main result of Murre's talk, one can prove that the latter morphism
is representable by unramified separated morphisms locally of finite
presentation; as a consequence, if $\underline{\operatorname{Hom}}_S(Z,Y)$ is
representable, so is $\underline{\operatorname{Hom}}_S(X,Y;\phi)$.

To get, given $X$ and $Y$, a representability theorem for
$\underline{\operatorname{Hom}}_S(X,Y)$, \struck{onge} one tries to find
morphisms $\phi_i : Z_i \to X$ as above, such that the open subfunctors
$\underline{\operatorname{Hom}}_S(X,Y;\phi_i)$ cover
$\underline{\operatorname{Hom}}_S(X,Y)$ (as a fpqc sheaf), and such that the
functors $\underline{\operatorname{Hom}}_S(Z_i,Y)$ are all representable. If for
instance $S$ is the spectrum of a field $k$, and if $X$ has ``enough'' points
radicial over $k$ (which is always true if $k$ is alg.\ closed) then we can take
for $Z_i$ all finite subschemes of $X$ whose points are radicial over $k$, and
we get

\page{31}
that $\underline{\operatorname{Hom}}_S(X,Y)$ is representable (any $Y$ locally
of finite presentation \add{(and separated)} over $k$); if we do not make any
assumption on $X$ except properness over $k$, the previous assumption becomes
true after finite ground-field extension $k'/k$, so that we get that for every
$Y$ as above, $\underline{\operatorname{Hom}}_S(X,Y) \times_S \operatorname{Spec}(k')$
is representable. From this Matsumura's theorem stated at the beginning follows
in a standart way by descebt arguments. The result holds too for
$\underline{\operatorname{Isom}}_k(X,Y)$ instead of
$\underline{\operatorname{Aut}}_k(X)$, but as you probably know,
$\underline{\operatorname{Hom}}_S(X,Y)$ is not always representable, even if $X$
is a quadratic extension of $S = \operatorname{Spec} k$, $Y$ being proper non
projective.

Over an arbitrary base $S$, one can give a fairly general statement of a
representability theorem, the points radicial over $k$ used above being
replaced by suitable flat subschemes of $X$. As \struck{applications}
\add{particular cases}, we get for instance that if $X$ has integral geometric
fibers and a section \add{along which $X$ is smooth,} then
$\underline{\operatorname{Hom}}_S(X,Y)$ is representable; and if $X$ has reduced
geometric fibers, then $\underline{\operatorname{Hom}}_S(X,Y)$ is representable
locally for the étale topology over $S$. Also, if $Y$ is quasi-projective over
$S = \operatorname{Spec} k$, then $\underline{\operatorname{Hom}}_S(X,Y)$ is
representable.

To fix the ideas, I gave the statements for
$\underline{\operatorname{Hom}}_S(X,Y)$, but one has quite analoguous resultst
of course for the $\prod_{X/S} P/X$ fonctors, which I guess will imply rather
formally the other ones.

If you are interested, I can send you a photocopy of the statement of the
general theorem of representability I alluded to above, and a couple of
corollaries (I already listed here the most striking ones).

Sincerely your's
\note{« (and separated) » est ajouté à l'encre ; « particular cases » et
« along which $X$ is smooth, » sont tapés dans l'interligne et reliés au texte
par un trait.}

\page{33}
\section*{Schémas Hom \add{etc} \struck{et \ill{} \ill{}}}

\begin{enumerate}
\item[1.] Formalisme des $\underline{\operatorname{Hom}}$ et $\prod_{S'/S}$
\item[2.] \struck{\ill{} \ill{} de Greenberg} \struck{\add{\uncertain{Généralités} \ill{}}}
Foncteurs de Serre-Lang \add{et} de Greenberg.
\uncertain{Extension} \ill{} $\underline{W}^{\cdot}(X)$ :
\item[3.] \uncertain{Cas} \ill{} \ill{} $\underline{W}^{\circ}(X)$ \uncertain{Th.}
\item[4.] Existence : Résultats Th.
\item[5.] Application aux immersions, morphismes de type fini,
\ill{} \uncertain{Atiyah}\uncertain{(?)}
\item[6.] Cas des morphismes \uncertain{quasi-projectifs} Th.
\item[9.] \struck{Ap} Cas de $\prod_{S'/S}$, \ill{} $S' \to S$ étale.
\marginal{\uncertain{reporter} \uncertain{à} \uncertain{la} \uncertain{fin}}
\item[7.] Cas des morphismes simples, et des morphismes étales. Th.
\item[{[8]}] Application aux morphismes $S' \to S$ quasi-étales
[\uncertain{i.e.}\ \ill{} \uncertain{loc.} \ill{}]. Th.
\item[{[10]}] Compléments dans le cas de Greenberg
[\uncertain{platitude} \ill{} \ill{} $\underline{W}$].
\end{enumerate}
\note{plan d'un exposé. Une flèche part du « 9 » et descend jusqu'après le [10] :
la rubrique 9 est sans doute à reporter en fin de plan. Le « 7 » est récrit
sur un autre chiffre. « Th. » en fin de ligne semble marquer les rubriques
qui comportent un théorème.}

\page{34}
\begin{tikzcd}
 & X \arrow[d] \arrow[dl, dashed, "\varphi"'] & Y \arrow[dl] & \\
S & T \arrow[l, "f"'] & T \times_S T \arrow[l, bend right=15, "\mathrm{pr}_1"'] \arrow[l, bend left=15, "\mathrm{pr}_2"] \arrow[ll, bend left=40, "g"] &
\end{tikzcd}

\note{dans l'original la flèche $X \to T$ et la flèche pointillée $\varphi : X \to S$
partent toutes deux de $X$ ; $Y$ descend sur $T$ ; $g : T\times_S T \to S$ est
un arc passant sous $T$.}

\begin{align*}
f_*(X/T) &= \underline{\operatorname{Hom}}_{T/S}(T/S, X/S) \\
f_*(X/T) &= \underline{\operatorname{Hom}}_S(T,X) \times_{\underline{\operatorname{Hom}}_S(T,T)} S \\
\underline{\operatorname{Hom}}_{T/S}(X/T, Y/T) &= f_*\bigl(\underline{\operatorname{Hom}}_T(X,Y)\bigr) \\
\underline{\operatorname{Hom}}_S(X,Y) &= \varphi_*(X \times Y / X) \\
\underline{\operatorname{Hom}}_S(X,Y) &= \underline{\operatorname{Hom}}_{S/S}(X/S, Y/S)
\end{align*}

Si $X$ est \uncertain{muni} d'une donnée de descente : $S$
[et si \uncertain{on} \uncertain{veut}, les conditions d'\uncertain{unicité}\,…],
et si $f_*(X)$ existe, alors l'objet descendu $X_0$ existe aussi, et on aura
\[
X_0 = f_*(X) \times_{g_*(X \times_S T)} f_*(X) .
\]

Contre-exemple de \uncertain{Nagata} : schéma algébrique simple $X$
\add{\uncertain{distinctes}} sur un corps alg.\ clos $k$, avec groupe $G \simeq \mathbb{Z}/2$
opérant sans point fixe, tel que $X/G$ n'existe pas. Soit $L$ une extension
de $k$ sur laquelle $G$ opère \uncertain{fidèlement},
\struck{pt fixe} $K = L^G$, faisant opérer $G$ diagonalement sur
$X \otimes_k L / L$, \ill{} \ill{} \ill{} sur $X \otimes_k L/L$ une
donnée de descente \uncertain{relativement} par $L/K$. Je
dis que la \uncertain{question} n'existe pas, i.e.\ la donnée
n'est pas \uncertain{effective}. En effet, \struck{si le quotient existait}
il \struck{existerait un \ill{} \ill{} \uncertain{conservation} des \ill{}}
\add{de $K$} \ill{} $L$,
\[
\text{\struck{$(X\otimes_k L)\otimes_K L = X\otimes_k(L\otimes_K L)$}} = X\otimes_k L
\]
\note{fin de page très surchargée : ratures et traits de renvoi entremêlés ;
la suite passe à la page 35.}

\page{35}
il y a \struck{\uncertain{deux} \ill{} de} une \uncertain{trajectoire} $\Gamma$
\uncertain{fermée} \uncertain{de} $G$ dans $X$ qui \uncertain{n'est} pas
\uncertain{contenue} dans \uncertain{un} ouvert affine. \uncertain{Les} \uncertain{ouverts}
\uncertain{affines} \uncertain{invariants} \uncertain{de} $X$ \uncertain{contenant}
\uncertain{une} \uncertain{trajectoire} \uncertain{des}
\marginal{\uncertain{distinctes} \uncertain{ou} \uncertain{deux} \uncertain{à} \uncertain{deux}
\uncertain{\ill{}}}
\ill{} $X \otimes_k L = X_L$ \uncertain{formeraient} \uncertain{une}
\struck{\ill{}} \add{\ill{}} \uncertain{famille} \ill{} \ill{}
\uncertain{n'ont} \uncertain{aucun} \ill{} \ill{} \ill{}. \uncertain{Ce} \uncertain{n'est}
\uncertain{une} \uncertain{trajectoire} \uncertain{sous} $G \times 1$.
De plus, \uncertain{un} \ill{} \ill{} $L \times G$ \uncertain{opérant}
\ill{} \ill{}, \uncertain{ses} \uncertain{opérations} \struck{\ill{}}
\uncertain{diagonales} \uncertain{de} $G$ \ill{} \ill{} \ill{} \ill{} $G \times 1$.

\textbf{Corollaire} $\underline{\operatorname{Hom}}_{\struck{\ill{}}/K}(\operatorname{Spec}(L), X_{(L)})$
\uncertain{n'existe} pas. A fortiori $\underline{\operatorname{Hom}}_K(\operatorname{Spec}(L), X_{(L)})$
n'existe \uncertain{pas}.

On \uncertain{peut} \uncertain{étendre} \uncertain{par} \uncertain{ceci} \uncertain{cet}
\uncertain{exemple} \ill{} $K$ \ill{}.
\note{page rapide : les formules et le corollaire se lisent, la prose qui les
relie à peine ; le bas de la page est blanc.}

\page{36}
\note{nouveau plan, numéroté par paragraphes. En haut à gauche, un cadre barré
de hachures croisées contient « $\dim Y/X = d$, $\dim X/S = \delta$ » et, en
bas, « $x_1 \ldots x_\delta$, $y_1 \ldots y_d$ » ; la marge gauche porte, en
travers, des notes serrées qu'on ne lit que par fragments.}

\begin{itemize}
\item[§ 1] Descente \add{d'objets plats} par morphismes quelconques
\item[§ 2] Descente d'objets quelconques par morphismes plats.
\marginal{\uncertain{intervertir} ?} \note{une accolade réunit § 2 et § 3.}
\item[§ 3] Application \add{du § 1} aux morphismes non \uncertain{ramifiés}
\item[§ 4] Cas particuliers élémentaires : Espaces $X/G$,
descente galoisienne non \uncertain{ramifiée}, le contre-exemple de \uncertain{Nagata}.
Interprétation de la descente plate finie par opérateurs différentiels. Cas de Cartier.
\item[§ 5] $\underline{\operatorname{Hom}}_S(X,Y)$ et $f_*(X)$ : formalisme général, cas
d'\uncertain{existence} \uncertain{pour} $X$ fini et plat sur $S$.
[Est lié : la descente par \struck{\ill{}} \uncertain{deux} \uncertain{conditions}
a) Existence de $\underline{\operatorname{Hom}}$ si $X \to S$ est \uncertain{fini} \uncertain{localement}
\uncertain{libre} \uncertain{?}
\ill{} fini et plat ---
b) Relations entre l'existence de $f_*(X)$, $\underline{\operatorname{Hom}}_S(X,Y)$,
et la possibilité de descente.]
\uncertain{Notations} : $\check{\mathbb{V}}(\mathcal{F})$ etc.
\item[§ 6] Application aux schémas $\underline{\operatorname{Hom}}$ : la géométrie
différentielle, les \uncertain{Picard}, Greenberg, l'algèbre commutative etc.
\item[§ 7] Hypo\struck{\ill{}}thèses sur $S$. Critère de représentabilité.
[y \uncertain{inclus} \uncertain{aussi} le critère général d'existence :
\uncertain{par} \uncertain{théorie} \uncertain{des} modules formels]
\item[§ 8] \struck{(\ill{})} \add{Théorème} \uncertain{fondamental} d'existence
A) Théorie formelle des modules. Extension des \uncertain{solutions}\,…
B) $f_*(X)$ si $X$ \uncertain{projectif} sur $Y$, $Y$ plat et propre sur $S$.
C) Picard.
\item[§ 9] [Critères de simplicité, de propreté, de platitude ?]
\item[§ 10] \note{numéro sans intitulé.}
\end{itemize}
\note{« § 8 » est récrit sur un « 9 » ; une accolade réunit § 8 à § 10, avec en
marge « \uncertain{ensemble} ? ». Les numéros § 9 et § 10 sont alignés à gauche
de B) et C) ; leur rattachement exact à ces lignes n'est pas sûr.}

\page{37}
\textbf{Proposition} Soit $f : T \to S$ un morphisme fini localement libre, et
soit \struck{\ill{}} $\mathcal{F}$ un faisceau \add{qu.\ coh.}\ localement libre sur $T$,
de sorte que $f_*(\mathcal{F})$ \uncertain{est} un faisceau \add{qu.\ coh.}\
localement libre sur $S$. Alors, on a un isomorphisme fonctoriel
\[
\boxed{\, f_*\bigl(\check{\underline{\mathbb{V}}}(\mathcal{F})\bigr) = \check{\underline{\mathbb{V}}}\bigl(f_*(\mathcal{F})\bigr) \,}
\]
[cet isomorphisme est compatible avec \uncertain{les} \uncertain{opérations} \uncertain{de}
\add{\uncertain{restriction}} image inverse, \add{localisation}].
\note{le reste de la page est blanc.}

\page{38}
\begin{tikzcd}
X \arrow[d, "f"] \\
Y = \operatorname{Spec}(\mathcal{O})
\end{tikzcd}

$f$ \add{\uncertain{plat} \uncertain{et}} propre, $Y$ \uncertain{schéma}
\uncertain{noethérien} local, $F$ cohérent sur $Y$, $F = \widetilde{M}$,
$F' = f^*(F)$, $F'_1 \subset F'$ \uncertain{un} \uncertain{sous-faisceau} \uncertain{cohérent}.
\struck{\ill{}} Je dis qu'il y a un plus petit sous-faisceau \add{\uncertain{cohérent}}
$F_1$ de $F$ tel que $f^*(F_1) \supset F'_1$.
\note{une flèche relie « $F'_1 \subset F'$ » à « sous-faisceau cohérent ».}

\uncertain{Mais} si $f^*(F_1)$ et $f^*(F_2) \supset F'_1$, alors
$f^*(F_1 \cap F_2) = f^*(F_1) \cap f^*(F_2) \supset F'_1$
\marginal{platitude de $f$} \note{« platitude de $f$ » est écrit sous l'égalité,
reliée à elle par une accolade.}
\uncertain{Les} $F_i \subset F$ tels que $f^*(F_i) \supset F'_1$ \struck{\ill{}} forment un
\uncertain{ensemble} \uncertain{filtrant} décroissant. \struck{\ill{}} \uncertain{Il}
\uncertain{suffit} \uncertain{de} \uncertain{prouver} que si $F_1 = \bigcap F_i$, alors
$f^*(F_1) \supset F'_1$, i.e.\ l'image de $F'_1$ dans
$F' / f^*(F_1) = f^*(F/F_1)$ est nulle. Remplaçant $F$ par $F/F_1$,
\uncertain{on} \uncertain{peut} \uncertain{supposer} \uncertain{que} $\bigcap F_i = 0$, alors
\uncertain{que} \uncertain{les} \uncertain{faisceaux} \struck{\ill{}} \add{images} de $F'_1$
\uncertain{contenus} \uncertain{dans} $\bigcap f^*(F_i)$ \uncertain{sont} \uncertain{nuls}
[\uncertain{est-il} \uncertain{vrai} \uncertain{que} $\bigcap f^*(F_i) = 0$ ?]
\struck{$\bigcap f^*(F_i) = 0$} ($f$ propre \struck{\ill{}}, les $F_i$ \uncertain{forment}
\uncertain{décroissant} \uncertain{dans} $F$ \uncertain{cohérent}, $Y$ local et noethérien).
\struck{\uncertain{Comme} $f$ \ill{}, \uncertain{les} $F'_1 = \bigcap f^*(\ill{})$ de}
\struck{\uncertain{Plat} $S$ \uncertain{une} \uncertain{section} \ill{}}

$F'_1$ \uncertain{est} \uncertain{fermé}, il suffit de prouver \uncertain{qu'il} \uncertain{est}
\uncertain{nul} sur $f^{-1}(y)$ ($y$ = point fermé de $Y$) donc de prouver
\uncertain{que} \uncertain{si} $s$ \uncertain{est} \uncertain{une} section de $F'_1$ \uncertain{sur}
\uncertain{un} \uncertain{ouvert} $U$, \ill{} \ill{} \ill{} \ill{} (\ill{} $x \in f^{-1}(y)$).
$\mathcal{O}_x$ \ill{} \ill{} \ill{} $\bigcap M_i = 0$, \uncertain{en} $M_i$
\uncertain{tendant} \uncertain{vers} $0$ dans $M$ (?)

\note{la marge gauche porte, en travers et dans un cadre arrondi, une note
« N.B. » : « Faut-il \uncertain{supposer} $M$, \ill{} \ill{} \ill{} $M_i \subset M$,
\ill{} \ill{} \ill{}, $\bigcap M_i = $ \ill{} \ill{} \ill{} … \ill{} [\uncertain{il}
\uncertain{suffit} \ill{} \ill{}] $f^*(\bigcap F_i) = \bigcap f^*(F_i)$ ? » ; dessous,
« \uncertain{Avec} \uncertain{plus} \uncertain{de} \uncertain{conditions} » et une flèche.
La lecture n'en est que fragmentaire.}

\page{39}
\begin{tikzcd}[column sep=small, row sep=small]
 & X \arrow[d] & \\
S & T \arrow[l, "f"'] & X \times_S S' \arrow[ul] \arrow[dl] \\
S' \arrow[u] & S' \times_S T = T' \arrow[l, "f'"] \arrow[u] \arrow[ur, bend right=20] &
\end{tikzcd}

\begin{tikzcd}[column sep=small, row sep=small]
 & X_s \arrow[d] & \\
s & T_s \arrow[l, no head] & \\
 & & X'_{s'} \arrow[uul] \arrow[dl] \\
s' & T'_{s'} \arrow[l, no head] \arrow[uu] &
\end{tikzcd}

\note{deux schémas en marge gauche, redessinés ici sans garantie sur toutes les
pointes ; dans le premier, une flèche $S' \to S$ et l'arc $T' \to X\times_S S'$
sont surchargés.}

\textbf{Proposition 2} Supposons que $f_*(X/T)$ existe, et que pour tout
$s \in S$, \struck{\ill{}} \uncertain{toute} \uncertain{réalisation} de $X_s/T_s$
\add{sur} \uncertain{définie} \uncertain{sur} \uncertain{une} \uncertain{extension} finie
[quelconque] de $k(s)$, soit \uncertain{« contenue »} dans un \struck{\ill{}}
$U_i$ \uncertain{des} $U_i$ ($U_i$, famille \ill{} \ill{} \ill{}). Alors
$f_*(X/T) \to$ \uncertain{réunion} des $f_{i*}(U_i/T)$ ($f_i = f|U_i$)
\marginal{\uncertain{dans} \uncertain{le} \uncertain{cas} \uncertain{« fini »}}
[\uncertain{on} \uncertain{suppose} $T$ et $X$ de type fini sur $S$ ?]

En effet, \struck{\uncertain{soit}} il suffit de montrer que \struck{\ill{}}
\uncertain{tout} \struck{\ill{}} \add{point} de $f_*(X/T)$ \uncertain{est} \uncertain{dans}
\uncertain{un} des $f_{i*}(U_i/T)$, \uncertain{et} il suffit \uncertain{de} \uncertain{le}
\uncertain{voir} sur un pt fermé \ill{} \ill{} \ill{} \uncertain{défini} \uncertain{sur}
\uncertain{un} corps $K$ sur un $k(s)$, et \uncertain{une} \add{\uncertain{réalisation}}
\uncertain{section} \uncertain{définie} \uncertain{sur} $K$ de $X_s$ sur $T_s$. \uncertain{On}
\uncertain{peut} \struck{\uncertain{définie} \uncertain{sur} \ill{}} $K$ fini sur $k(s)$,
vu les hyp.\ de finitude \uncertain{précédentes}. \ill{} \ill{} \ill{} \ill{}
\uncertain{comme} \uncertain{des} hyp.\ \ill{} \ill{} cas.

\struck{\ill{}} \uncertain{Il} \uncertain{provient} \uncertain{donc} \ill{} \ill{}
\uncertain{par} \uncertain{les} \uncertain{réalisations} \add{\uncertain{multi}}.

\struck{\uncertain{Corollaire}} \add{Proposition 2 bis}
\marginal{\uncertain{remplace} prop.~2}
Supposons \uncertain{que} \uncertain{pour} \uncertain{toute} \uncertain{extension} finie (\uncertain{comp.}…)
$X_s/T_s$ \uncertain{admette} \uncertain{une} \ill{} de $k(s)$ \uncertain{soit} \uncertain{contenue} \uncertain{dans},
(\uncertain{quelconque}) $K'$ \uncertain{en} \ill{}
\uncertain{un} \uncertain{ouvert} \struck{\uncertain{affine} \ill{}} \add{d'}une famille $U_i$.
Supposons \uncertain{que} \uncertain{les} $f_{i*}(U_i/T)$ \uncertain{existent}, alors
il \uncertain{en} \uncertain{est} \uncertain{de} \uncertain{même} \uncertain{de} $f_*(X/T)$.
\note{page rapide ; la proposition 2 bis est le titre « Corollaire »
biffé deux fois et remplacé en interligne.}

\page{40}
\textbf{Démonstration.} Les \uncertain{conditions} \uncertain{suffisantes} sont
\uncertain{les} \uncertain{mêmes}, \uncertain{si} les $f_{i*}(U_i/T)$ existent. \uncertain{Alors},
il s'ensuit que \ill{} \ill{} \uncertain{les} $f_{i*}(U_i) \to f_*(U_j)$.
\uncertain{Les} conditions de \uncertain{recollement} \uncertain{sont} \uncertain{satisfaites}
\uncertain{vérifiées}, \uncertain{de} \uncertain{là} \uncertain{un} \uncertain{préschéma} $\underline{H}$ sur $S$,
\struck{\ill{}} \uncertain{réunion} \uncertain{des} \ill{} $H_i \simeq f_*(U_i/T)$.
Les \uncertain{définit} \uncertain{par} \uncertain{les} \uncertain{morphismes} \uncertain{universels} $T$-\uncertain{linéaires}
\struck{\uncertain{universels}} $H_i \times T \to U_i$,
\[
H \times_S T \longrightarrow X ,
\]
\uncertain{je} dis que \uncertain{ce} \uncertain{dernier} \uncertain{est} \uncertain{universel}, \uncertain{i.e.}\
\struck{\ill{}} $H$ \uncertain{représente} $f_*(X/T)$.
\marginal{\uncertain{Par} \uncertain{définition}, \uncertain{les} \uncertain{morphismes} $H_i \to H$
\uncertain{et} \uncertain{les} \ill{} \ill{} \ill{} $H \times_S T$ \ill{} \ill{} \ill{} \ill{}
\uncertain{Soit} \ill{} \ill{} \ill{}}
\struck{\uncertain{Pour} \uncertain{cela}, \uncertain{soit} $S' \to S$ \ill{} \ill{} \ill{}, \ill{} \uncertain{faut} \uncertain{prouver}}
\uncertain{On} \uncertain{doit} \uncertain{montrer}
\[
\underline{\operatorname{Hom}}_S(S', H) \xrightarrow{\ \sim\ } \underline{\operatorname{Hom}}_T(S' \times_S T, X) .
\]
\uncertain{On} \uncertain{peut} \uncertain{supposer} $S = S'$. \uncertain{Les} \uncertain{deux} \ill{}
\uncertain{sont} \uncertain{des} \uncertain{faisceaux} \uncertain{pour} \uncertain{la} \uncertain{topologie}
\uncertain{de} \uncertain{Zariski} \struck{\ill{} \ill{}},
\uncertain{il} \uncertain{suffit} \uncertain{de} \uncertain{prouver} \uncertain{que}
\uncertain{les} \uncertain{deux} \uncertain{ont} \uncertain{mêmes} \uncertain{sections} :
\struck{\ill{}} \uncertain{bijection}, \uncertain{on} \uncertain{se} \uncertain{ramène} \uncertain{au}
\uncertain{cas} \uncertain{où} \struck{\ill{}} :
\struck{$S' \to S$} $u : S \to H$ \add{\uncertain{correspondant}}
\struck{\uncertain{et}} $T \to X$, \uncertain{alors}
\struck{$S \to T$} \uncertain{définissant} \uncertain{le} $T \to X$,
\uncertain{les} \uncertain{images} \uncertain{des} $T_i$ \uncertain{par} \ill{} \uncertain{des}
\struck{\uncertain{ouverts}} \uncertain{des} $U_i$ \uncertain{recouvrent} $T$,
i.e.\ il \uncertain{faut} \uncertain{prouver}
\[
\underline{\operatorname{Hom}}_S(S, H) \simeq \underline{\operatorname{Hom}}_T(T, X) ,
\]
\uncertain{soit} $\varphi$ \uncertain{une} \uncertain{section} de $X/T$, \ill{}

\begin{tikzcd}[column sep=small, row sep=small]
 & X \arrow[d] & \\
S & T \arrow[l] & \\
H \arrow[u] & H \times_S T = \bigcup H_i \times_S T \arrow[l] \arrow[u] \arrow[uu, bend right=30] &
\end{tikzcd}

\note{page très surchargée : deux longs traits de renvoi relient les ratures du
bas au texte du haut ; un petit croquis (lignes verticales, entre $S$ et
$H_i$) est barré dans la marge gauche, qui porte aussi en travers une note
illisible (« \uncertain{par} \uncertain{définition} … $H_i$ … $H \times_S T$ … »).
La démonstration se poursuit au-delà du lot.}

\end{document}
